\documentclass[../../../uem_phy_std_a.tex]{subfiles}

\begin{document}
    \begin{enumerate}
        \item キルヒホッフ則より，
            
            \myspaceb%
            $\displaystyle%
                V_0 = \frac{q_0}{3C} \qquad \therefore \; q_0 = \uwave{3CV_0} \;.
            $
        \item 電荷保存則より，
            
            \myspaceb%
            $\displaystyle%
                x_1 + y_1 = q_0 + 0 \qquad \therefore \; x_1 + y_1 = 3CV_0 \;.
            $\\            
            キルヒホッフ則より，
            
            \myspaceb%
            $\displaystyle%
                \myfracc{x_1}{3C} = \myfracc{y_1}{2C}%
                \qquad \therefore \; x_1 = \uwave{\myfraca{9}{5} CV_0} \;,%
                \quad y_1 = \uwave{\myfraca{6}{5} CV_0} \;.
            $
        \item S$_2$を開き，S$_1$を閉じて充分待つと，C$_1$の電荷は再び$q_0=3CV_0$となる%
            （C$_2$の電荷は$y_1$のまま）．続いて，S$_1$を開き，S$_2$を閉じて充分待った後，%
            電荷保存則より，
            
            \myspaceb%
            $\displaystyle%
                x_2 + y_2 = q_0 + y_1 \qquad \therefore \; x_2 + y_2 = \myfraca{21}{5}CV_0 \;.
            $\\
            キルヒホッフ則より，
            
            \myspaceb%
            $\displaystyle%
                \myfracc{x_2}{3C} = \myfracc{y_2}{2C}%
                \qquad \therefore \; x_2 = \uwave{\myfraca{63}{25} CV_0} \;,%
                \quad y_2 = \uwave{\myfraca{42}{25} CV_0} \;.
            $
        \item キルヒホッフ則より，
            
            \myspaceb%
            $\displaystyle%
                \myfraca{X}{3C} = \myfraca{Y}{2C}%
                \qquad \therefore \; X = \uwave{3CV_0} \;,%
                \quad Y = \uwave{2CV_0} \;.
            $
    \end{enumerate}
\end{document}
